% ============================================================
%  SCHOLA DOMESTICA — Level 1, Week 3, Lesson 11
%  Topic: The Number Zero — None at All
%  8–10 problems | No speed drill | Picture-heavy
%  Compile with XeLaTeX
% ============================================================
\documentclass[12pt,letterpaper]{article}
\usepackage{sd_style}

\renewcommand{\lessonnum}{11}
\renewcommand{\lessonweek}{3}

\begin{document}

\lessonheader{3}{11}{The Number Zero}

\begin{instructbox}
\noindent\textbf{\color{sd-blue}Instruction} \textit{(read aloud):}\\[4pt]
Today we learn the number \textbf{zero}.
Zero means \textit{none} — an empty group, nothing at all.
Zero is written as \textbf{0}, and it comes before 1 on the number line.
Even though zero means nothing, it is a very important number!
\end{instructbox}

\vspace{8pt}
\noindent\textbf{\color{sd-blue}Trace and write:}
\tracerowone{0}

\goldrule
\noindent\textbf{\color{sd-blue}Practice}
\vspace{6pt}

\prob How many apples are in the empty bowl?\\[6pt]
\hspace{1cm}%
\IfFileExists{empty_bowl.jpg}{\includegraphics[height=1.8cm]{empty_bowl.jpg}}{%
  \begin{tikzpicture}[baseline=0]
    \draw[sd-blue, thick] (0,0) arc (180:360:1.2cm and 0.5cm);
    \draw[sd-blue, thick] (-1.2,0) arc (180:360:1.2cm and 1.0cm);
  \end{tikzpicture}}
\quad\ansblank

\vspace{10pt}
\prob How many ducks are on the \textit{empty} pond?\\[4pt]
\hspace{1cm}\IfFileExists{empty_pond.jpg}{\includegraphics[height=1.8cm]{empty_pond.jpg}}{%
  \textit{[Teacher: show an empty space or blank box]}}
\quad\ansblank

\vspace{10pt}
\prob Circle the empty group.\\[6pt]
\begin{center}
\begin{tabular}{ccc}
\textbf{A}\quad\onedot\onedot\onedot &\quad\quad&
\textbf{B}\quad\fbox{\rule{0pt}{0.7cm}\hspace{1.4cm}}
\end{tabular}
\end{center}

\vspace{10pt}
\prob Write the number that means \textit{none}: \quad\ansblank

\vspace{10pt}
\prob What number comes just before 1 on the number line?\\[6pt]
\begin{center}
\begin{tikzpicture}
  \draw[sd-blue, very thick, ->] (-0.3,0) -- (5.5,0);
  \foreach \n in {0,...,5}{
    \draw[sd-blue, thick] (\n, -0.15) -- (\n, 0.15);
    \node[below] at (\n, -0.25) {\textbf{\n}};
  }
  \draw[sd-gold, very thick, ->] (0, 0.5) -- (0, 0.18);
  \node[above] at (0, 0.55) {\color{sd-gold}\textbf{?}};
\end{tikzpicture}
\end{center}
\ansblank

\vspace{10pt}
\prob \textit{(Teacher reads)} There were 3 birds on a branch.
All 3 flew away. How many birds are left? \quad\ansblank

\vspace{10pt}
\prob Is zero more or fewer than 1? \quad\ansblank

\vspace{10pt}
\prob Put these numbers in order, least to greatest:\\[6pt]
\hspace{2cm}{\Large 3 \quad 0 \quad 1 \quad 2}\\[6pt]
\hspace{2cm}\ansblank\quad\ansblank\quad\ansblank\quad\ansblank

\vspace{10pt}
\prob Draw \textbf{zero} butterflies in the box.\\[4pt]
\noindent\fbox{\rule{0pt}{2cm}\hspace{11cm}}\\[4pt]
\textit{(What does your drawing look like? That is zero!)}

\goldrule

\begin{checkbox}
\noindent\textbf{\color{sd-green}Knowledge Check}\\[4pt]
\setcounter{probnum}{0}
\prob Write the number zero: \quad\ansblank\\[8pt]
\prob Which is less: 0 or 5? \quad\ansblank\\[8pt]
\prob \textit{(Teacher holds up a closed fist)} How many fingers am I showing? \quad\ansblank
\end{checkbox}

\vfill
\noindent{\color{sd-blue}$\bigstar$ \textbf{Enrichment:}}
The ancient Romans had \textit{no symbol for zero}!
They simply wrote nothing.
The number zero was invented in India and brought to Europe by Arab mathematicians.
In Latin, we say \textit{nihil} (nothing) or \textit{nulla} (none).

\vspace{4pt}\hrule\vspace{3pt}
{\footnotesize\color{gray}\textit{Teacher note: Zero is conceptually tricky for young children.
  Reinforce with physical objects — show a full cup, then pour everything out.
  "Now how many are in the cup? Zero — none at all."}}

\end{document}
